Un depósito de sección horizontal constante tiene altura total $3\,\mathrm m$ y un orificio circular de radio $0.05\,\mathrm m$ en el fondo. La velocidad de salida es $c\sqrt{2gh}$, donde $h$ es la altura de agua, $0<c<1$ y $g=9.8\,\mathrm{m/s^2}$. Determine la EDO para $h$, denotando por $A_w$ el área horizontal del depósito.

Se supone líquido incompresible, superficie horizontal, sección constante $A_w>0$, orificio pequeño y velocidad de salida suministrada por la ley indicada. Se aplica el modelo mientras queda agua y no se excede la altura del depósito.
